\documentclass[11pt,a4paper]{article}
\usepackage[margin=22mm]{geometry}
\usepackage{amsmath,amssymb}
\usepackage{fontspec}
\usepackage{polyglossia}
\setdefaultlanguage{russian}
\setmainfont{DejaVu Serif}
\usepackage[unicode,hidelinks]{hyperref}

\allowdisplaybreaks[1]
\raggedright
\setlength{\parskip}{0.35em}
\newcommand{\dd}{\mathop{}\!\mathrm{d}}
\newcommand{\avg}[1]{\langle #1\rangle}
\newcommand{\ket}[1]{\lvert #1\rangle}
\newcommand{\bra}[1]{\langle #1\rvert}
\newcommand{\braket}[2]{\langle #1\mid #2\rangle}
\DeclareMathOperator{\Var}{Var}

\newcommand{\solutiontitle}[4]{%
  \hypersetup{pdftitle={Квантовая механика 1. Решения ДЗ №#1}}%
  \begin{center}
    {\small Квантовая механика 1 \enspace\textperiodcentered\enspace Семестр 1\par}
    \vspace{1ex}
    {\LARGE\bfseries Решения домашнего задания №#1\par}
    \vspace{0.5ex}
    {\large #2\par}
  \end{center}
  \noindent{\small Условия приведены в файле «#3», #4.}\par\medskip
}

\geometry{hmargin=19mm,vmargin=18mm}
\setlength{\parskip}{0.2em}
\begin{document}
\solutiontitle{10}{Рассеяние на барьерах и линейный потенциал}{ДЗ2026(1)\_10 (2).pdf}{с.~1}

\section{Один дельта-барьер}

Пусть \(E>0,\ k=\sqrt{2mE}/\hbar,\ g=m\alpha/\hbar^2\).
Волна при падении слева имеет вид \(e^{ikx}+r e^{-ikx}\)
слева и \(t e^{ikx}\) справа. Непрерывность даёт \(1+r=t\).
После интегрирования стационарного уравнения через ноль
\(\Psi'(0^+)-\Psi'(0^-)=2g\Psi(0)\). Отсюда
\(ik(t-1+r)=2g t\), или \(ik r=g t\).
При \(\eta=g/k\) находим
\[
 t=\frac1{1+i\eta},\qquad r=\frac{-i\eta}{1+i\eta}.
\]
Поток каждой бегущей волны с амплитудой \(A\) имеет модуль
\(\hbar k|A|^2/m\), причём \(k\) одинаково с двух сторон.
Поэтому
\[
 T=|t|^2=\frac1{1+\eta^2},\qquad R=|r|^2=\frac{\eta^2}{1+\eta^2}.
\]
Сумма \(T+R=1\), а при \(\alpha=0\) имеем \(T=1\).

\section{Два одинаковых дельта-барьера}

У каждого барьера используем локальный отсчёт координаты
от его положения и амплитуды \(A,B\) перед \(e^{ikx},e^{-ikx}\)
в этом отсчёте. Непрерывность сохраняет \(A+B\), а скачок
производной меняет \(ik(A-B)\) на \(ik(A-B)+2g(A+B)\).
Значит матрица, переводящая столбец амплитуд слева в столбец
справа, есть
\[
 M=\begin{pmatrix}1-i\eta&-i\eta\\ i\eta&1+i\eta\end{pmatrix},
 \qquad \det M=1.
\]
Второй отсчёт сдвинут на \(a\), поэтому между барьерами амплитуды
получают множители \(e^{ika},e^{-ika}\). В результате
\[
 \binom{t e^{ika}}0=\mathcal M\binom1r,\qquad
 \mathcal M=M\operatorname{diag}(e^{ika},e^{-ika})M .
\]
Из второй строки \(r=-\mathcal M_{21}/\mathcal M_{22}\), из
первой и \(\det\mathcal M=1\) получаем \(t e^{ika}=1/\mathcal M_{22}\).
Перемножение даёт
\[
 \mathcal M_{22}=(1+i\eta)^2 e^{-ika}+\eta^2 e^{ika},
 \qquad t=\frac1{(1+i\eta)^2+\eta^2 e^{2ika}}.
\]
Возводим знаменатель по модулю в квадрат. Его вещественная часть
\(1-2\eta^2\sin^2(ka)\), мнимая \(2\eta+2\eta^2\sin(ka)\cos(ka)\).
В результате
\[
 T=\frac1{1+4\eta^2[\cos(ka)+\eta\sin(ka)]^2},\qquad R=1-T .
\]
По обе стороны от барьеров \(j=(\hbar k/m)(|A|^2-|B|^2)\).
Скачок производной пропорционален \(\Psi\) с вещественным
коэффициентом и сохраняет ток, что даёт \(R=1-T\).
При \(a\to0\) получается
\(T=1/(1+4\eta^2)\) для одного барьера удвоенной силы.
Для \(\eta\ne0\) при \(\cos(ka)+\eta\sin(ka)=0\) имеем резонанс
\(T=1\), тогда как один барьер обладает \(T<1\). В зависимости
от фазы прохождение через два барьера может быть больше
или меньше прохождения через один.

\section{Подбарьерное прохождение через прямоугольный барьер}

Пусть \(0<E<U_0\), \(k=\sqrt{2mE}/\hbar\),
\(\kappa=\sqrt{2m(U_0-E)}/\hbar\).
Внутри \(\Psi''=\kappa^2\Psi\). Для столбца \((\Psi,\Psi')^{\mathsf T}\)
это даёт матрицу перехода от \(0\) к \(a\)
\[
 \begin{pmatrix}c&s/\kappa\\ \kappa s&c\end{pmatrix},\qquad
 c=\cosh(\kappa a),\quad s=\sinh(\kappa a).
\]
На краях \(\Psi,\Psi'\) непрерывны.
Справа \((\Psi(a),\Psi'(a))=te^{ika}(1,ik)\).
Умножение на обратную матрицу, полученную заменой \(s\to-s\),
даёт
\(\Psi(0)=te^{ika}(c-iks/\kappa)\) и
\(\Psi'(0)=te^{ika}(-\kappa s+ikc)\).
Амплитуда падающей волны равна
\([\Psi(0)+\Psi'(0)/(ik)]/2=1\), поэтому
\[
 t=\frac{e^{-ika}}{c+i(\kappa^2-k^2)s/(2\kappa k)}.
\]
Учитывая \(c^2=1+s^2\), \(k^2+\kappa^2=2mU_0/\hbar^2\) и
сохранение тока, получаем
\[
 T=\left[1+\frac{U_0^2}{4E(U_0-E)}
                        \sinh^2(\kappa a)\right]^{-1},\qquad R=1-T.
\]
При фиксированном \(E\) и \(\kappa a\gg1\) имеем
\(\sinh(\kappa a)\sim e^{\kappa a}/2\), откуда
\[
 T(E)\sim\frac{16E(U_0-E)}{U_0^2}e^{-2\kappa a}.
\]
При \(a\to0\) формула даёт \(T\to1\).

\section{Линейный потенциал при больших отрицательных координатах}

Для \(\alpha>0\) положим \(x_t=E/\alpha\),
\(\ell=(\hbar^2/(2m\alpha))^{1/3}\), \(z=(x-x_t)/\ell\).
Уравнение \(-\hbar^2\Psi''/(2m)+\alpha x\Psi=E\Psi\)
становится \(\dd^2\Psi/\dd z^2=z\Psi\).
Общее решение \(A\operatorname{Ai}(z)+B\operatorname{Bi}(z)\).
На запрещённой стороне \(x\to+\infty\) условие убывания
оставляет \(B=0\), так как \(\operatorname{Bi}\) там растёт.

При \(x\to-\infty\), \(s=-z\to\infty\), асимптотики Эйри дают
\[
 \Psi(x)=\frac1{\sqrt\pi\,s^{1/4}}
     \left[A\sin\!\left(\tfrac23s^{3/2}+\tfrac\pi4\right)
          +B\cos\!\left(\tfrac23s^{3/2}+\tfrac\pi4\right)
          +O(s^{-3/2})\right].
\]
Ведущая фаза \(\tfrac23s^{3/2}\) равна
\(\hbar^{-1}\int_x^{x_t}\sqrt{2m(E-\alpha x')}\dd x'\).
Сдвиг \(\pi/4\) и амплитуду ветви \(\operatorname{Ai}\)
видно из интеграла Эйри
\(\operatorname{Ai}(-s)=\pi^{-1}\int_0^\infty\cos(v^3/3-sv)\dd v\).
Стационарная точка \(v=\sqrt s\) имеет фазу \(-2s^{3/2}/3\)
и положительную вторую производную \(2\sqrt s\). Гауссов
интеграл в методе стационарной фазы даёт
\(\pi^{-1/2}s^{-1/4}\cos(-2s^{3/2}/3+\pi/4)\), что равно
синусной асимптотике выше.

При \(\alpha<0\) используем
\(z=(x_t-x)/\ell,\ \ell=(\hbar^2/(2m|\alpha|))^{1/3}\).
Тогда \(z\to+\infty\) при \(x\to-\infty\), и решение,
убывающее на этом конце, пропорционально
\(z^{-1/4}e^{-2z^{3/2}/3}\). Независимая ветвь растёт как
\(z^{-1/4}e^{2z^{3/2}/3}\). При \(\alpha=0\) остаётся
свободная задача, у которой решениями служат плоские волны
при \(E>0\), линейные функции при \(E=0\) и экспоненты при \(E<0\).

\end{document}
